\documentclass[10pt,a4paper]{article}
\usepackage[margin=25mm]{geometry}
\usepackage{amsmath,amssymb,amsthm,hyperref}
\newtheorem{proposition}{Proposition}
\newtheorem{theorem}[proposition]{Theorem}
\newcommand{\core}{\operatorname{core}}
\newcommand{\Aut}{\operatorname{Aut}}
\title{Keeping subgroup identities in a coset action}
\author{Human Compression project}
\date{}
\begin{document}\maketitle

Algorithm 3.1 of Leemans and Rodrigues~\cite{LR} takes an arbitrary
group as input. It passes maximal subgroups to a coset action and then
takes their full inverse images. This operation can replace a maximal
subgroup by the whole group. We give the exact obstruction and a repair
that leaves the orbit computation unchanged. The issue does not arise
for the nonabelian simple groups in the source's numerical tables.

\section{The inverse image of a subgroup image}
Let $G$ be finite, let $H$ be a maximal proper subgroup, and let
$\varphi:G\longrightarrow\operatorname{Sym}(H\backslash G)$
be the action on right cosets by right multiplication. Its kernel is
\[
 N=\core_G(H)=\bigcap_{g\in G}g^{-1}Hg.
\]
In particular this action need not be faithful, contrary to the
unqualified assertion on source page~4.

\begin{proposition}\label{preimage}
For every subgroup $K\leq G$,
$\varphi^{-1}(\varphi(K))=KN$.
If $K$ is maximal proper, this inverse image equals $K$ when
$N\leq K$, and equals $G$ otherwise. For a point $x=Hg$, however,
\[
 \varphi^{-1}\bigl(\varphi(G)_x\bigr)=g^{-1}Hg.
 \tag{1}
\]
\end{proposition}
\begin{proof}
An element $y$ has image in $\varphi(K)$ exactly when
$\varphi(y)=\varphi(k)$ for some $k\in K$, or equivalently
$k^{-1}y\in N$. This proves the first equality. Normality of $N$
makes $KN$ a subgroup containing $K$; maximality gives the stated
dichotomy. Finally $\varphi(y)$ fixes $Hg$ exactly when
$Hgy=Hg$, equivalently $gyg^{-1}\in H$. This proves (1), with
no faithfulness assumption.
\end{proof}

The printed algorithm calls its image subgroup $M_2$ and outputs
$\{\varphi^{-1}(M_2),\varphi^{-1}(G_x)\}$. In this notation,
if the original subgroup was $K$, the actual output is
$\{KN,g^{-1}Hg\}$. The separate instruction to compute the orbits
of an undefined $K$ is an evident notation error; interpreting it as
the just-constructed image subgroup does not repair the inverse-image
problem.

\section{A counterexample of order four}
Take the additive group $G=\mathbb F_2^2$ with distinct order-two
subgroups $H=\langle(1,0)\rangle$ and
$K=\langle(0,1)\rangle$. Both are maximal, and $N=H$. On the two
cosets of $H$, the image of $K$ is transitive. Its one orbit therefore
has size two and passes the printed orbit-size test. For either
representative point the printed output is
\[
                   \{G,H\},
\]
since $KH=G$ and every point stabilizer in the image is trivial.
This is not a pair of maximal proper subgroups. The resulting coset
incidence graph has one vertex on its $G$ side and two on its $H$
side. Each of the latter has degree one, so the geometry is not firm.
By contrast the pair $\{K,H\}$ gives the complete bipartite graph
$K_{2,2}$ and is firm.

The small accompanying certificate enumerates the four elements and
their three maximal subgroups. Each of the six ordered choices of
distinct subgroups produces the same obstruction. It also constructs
the two incidence graphs just described. This finite calculation
checks the counterexample; Proposition~\ref{preimage} proves the
obstruction for all finite groups.

\section{The corrected algorithm}
Keep the list of original maximal subgroups together with their images.
For each original pair $H,K$ in the representative list, compute the
$\varphi(K)$-orbits on $H\backslash G$. For a representative $x$
of each orbit of size greater than one, output
\[
             \{K,\varphi^{-1}(\varphi(G)_x)\}.
 \tag{2}
\]
As in the source, retain only one pair from each orbit of $\Aut(G)$
on unordered pairs. This last step is a required equivalence test,
not something supplied by the coset orbit calculation alone.

\begin{theorem}\label{correct}
The corrected algorithm produces precisely the unordered pairs of
distinct maximal proper subgroups, up to the equivalence specified in
the source. Consequently it classifies the source's firm, connected,
flag-transitive primitive coset geometries of rank two for every finite
group, subject to a complete maximal-subgroup list and an exact
automorphism-equivalence test.
\end{theorem}
\begin{proof}
For $x=Hg$, put $H_x=g^{-1}Hg$. By (1) this is the second subgroup
in (2), so both output subgroups are maximal proper. The orbit size is
\[
                 |x^K|=[K:K\cap H_x].
\]
It equals one exactly when $K\leq H_x$, which by maximality is
equivalent to $K=H_x$. Thus the orbit-size condition retains exactly
the distinct pairs. Such a pair generates $G$: a subgroup containing
both would strictly contain a maximal subgroup. Both residue sizes
$[K:K\cap H_x]$ and $[H_x:K\cap H_x]$ are at least two, proving
firmness. Generation gives connectedness, and right multiplication
is transitive on incident pairs of cosets: any element in their
intersection carries the standard incident pair to that pair.

For completeness, start with any ordered pair $A,B$ of distinct
maximal subgroups. Conjugate it so that $B$ becomes a chosen
representative $K$. The resulting first subgroup is $g^{-1}Hg$ for
some chosen representative $H$ and some $g\in G$. The point $Hg$
lies in a $K$-orbit whose chosen representative is $Hgk$ for some
$k\in K$. Its stabilizer is $k^{-1}g^{-1}Hgk$. Conjugation by $k$
fixes $K$ and carries the previous pair to the produced pair. The
orbit is nonsingleton because the subgroups are distinct. Hence every
pair is produced up to inner conjugacy; the final exact test takes
the required quotient by all automorphisms and type interchange.
\end{proof}

If $G$ is nonabelian simple, every proper maximal subgroup has trivial
core. Proposition~\ref{preimage} then shows that the printed
inverse-image step already returns $K$. The correction above therefore
does not explain discrepancies in the sporadic numerical tables; those
are separate source claims.

The subgroup-orbit method, the orbit-size test and the final
automorphism test originate in~\cite{LR}. The kernel identity and the
orbit-stabilizer theorem are elementary group theory. The contribution
here is the diagnosis of the stated arbitrary-group scope, its explicit
counterexample, and the correction retaining the original subgroup.
No new classification of the sporadic incidence graphs, Lean proof or
compression measurement is claimed by this component.

\begin{thebibliography}{1}
\bibitem{LR} D. Leemans and B. G. Rodrigues,
\emph{On primitive geometries of rank two},
The Art of Discrete and Applied Mathematics 2 (2019), P1.08,
\href{https://doi.org/10.26493/2590-9770.1301.165}{doi:10.26493/2590-9770.1301.165},
pages 4--5, Algorithm 3.1.
\end{thebibliography}
\end{document}
